\documentclass[10pt]{article}
% Shared typographic and algebraic source.
\usepackage[letterpaper,margin=0.86in]{geometry}
\usepackage[T1]{fontenc}
\usepackage[utf8]{inputenc}
\usepackage{lmodern}
\usepackage{microtype}
\usepackage{amsmath,amssymb,amsthm,mathtools,mathrsfs}
\usepackage{booktabs,array,tabularx,longtable}
\usepackage{enumitem}
\usepackage[dvipsnames]{xcolor}
\usepackage{fancyhdr,titlesec}
\usepackage[hidelinks]{hyperref}

\definecolor{ink}{HTML}{1F2A33}
\definecolor{teal}{HTML}{226A73}
\definecolor{muted}{HTML}{65747E}
\definecolor{line}{HTML}{C5D1D3}
\definecolor{wash}{HTML}{F3F7F6}

\setlength{\parindent}{0pt}
\setlength{\parskip}{0.46em}
\setlength{\emergencystretch}{2em}
\allowdisplaybreaks[2]
\renewcommand{\arraystretch}{1.16}
\setlist{nosep,leftmargin=1.5em}

\pagestyle{fancy}
\fancyhf{}
\fancyhead[L]{\footnotesize\sffamily\color{muted}COPPALINI ARCHITECTURE}
\fancyhead[R]{\footnotesize\sffamily\color{muted}PHYSICAL REALITY INSIDE CONSTRAINT}
\fancyfoot[C]{\footnotesize\sffamily\color{muted}\thepage}
\renewcommand{\headrulewidth}{0.35pt}
\renewcommand{\headrule}{\hbox to\headwidth{\color{line}\leaders\hrule height \headrulewidth\hfill}}
\setlength{\headheight}{19pt}

\titleformat{\section}{\Large\bfseries\color{ink}}{\thesection.}{0.55em}{}
\titleformat{\subsection}{\large\bfseries\color{teal}}{\thesubsection.}{0.5em}{}
\titlespacing*{\section}{0pt}{1.42em}{0.48em}
\titlespacing*{\subsection}{0pt}{1.02em}{0.32em}

\newtheorem{theorem}{Theorem}
\newtheorem{proposition}[theorem]{Proposition}
\newtheorem{lemma}[theorem]{Lemma}
\theoremstyle{definition}
\newtheorem{definition}[theorem]{Definition}
\theoremstyle{remark}
\newtheorem{remark}[theorem]{Boundary}

\newcommand{\RR}{\mathbb R}
\newcommand{\CC}{\mathbb C}
\newcommand{\QQ}{\mathbb Q}
\newcommand{\ZZ}{\mathbb Z}
\newcommand{\Id}{I}
\newcommand{\tr}{\operatorname{tr}}
\newcommand{\im}{\operatorname{im}}
\newcommand{\kerop}{\operatorname{ker}}
\newcommand{\Span}{\operatorname{span}}
\newcommand{\diag}{\operatorname{diag}}
\newcommand{\snoc}{\operatorname{snoc}}
\newcommand{\Admit}{\operatorname{Admit}}
\newcommand{\Gate}{\operatorname{Gate}}
\newcommand{\Typed}{\operatorname{Typed}}
\newcommand{\Future}{\operatorname{future}}
\newcommand{\Scope}{\operatorname{sc}}
\newcommand{\SilScope}{\mathsf{Silence}_{\rm scope}}
\newcommand{\SilFit}{\mathsf{Silence}_{\rm fit}}

\newcommand{\papertitle}[2]{%
  \begin{center}
  {\color{teal}\rule{\textwidth}{1.35pt}}\\[0.38in]
  {\footnotesize\sffamily\bfseries\color{teal} COPPALINI ARCHITECTURE}\\[0.16in]
  {\fontsize{31}{36}\selectfont\bfseries\color{ink}#1\par}
  \vspace{0.15in}
  {\Large\color{teal}#2\par}
  \vspace{0.22in}
  {\color{line}\rule{\textwidth}{0.55pt}}\\[0.12in]
  {\normalsize Anthony Vito Coppa}\par
  {\small 9 November 2026}\par
  \vspace{0.30in}
  \end{center}
}

\newcommand{\sharedroot}{%
Let
\[
Q=\begin{pmatrix}1&1\\1&0\end{pmatrix},\qquad
\Omega=\begin{pmatrix}0&1\\-1&0\end{pmatrix},\qquad
C=-\Omega,
\]
\[
G=[Q,\Omega]=\begin{pmatrix}-2&1\\1&2\end{pmatrix},\qquad
H=\frac{G}{\sqrt5},\qquad
S=\frac{2Q-I}{\sqrt5}.
\]
The common exact identities are
\[
Q^2=Q+I,\qquad C^2=-I,\qquad G^2=5I,\qquad H^2=S^2=I,
\]
\[
H=CS,\qquad SG=\sqrt5\Omega=-GS,\qquad
L_s:=e^{sS},\qquad
L_s^{\mathsf T}GL_s=G,\qquad HL_sH=L_{-s}.
\]
}

\newcommand{\sharedoccurrence}[2]{%
A physical field has an interface
\[
\begin{aligned}
\mathfrak I_J:=(&J,\Gamma_J,\operatorname{Src}_J;
A_J,X_J,Y_J,T_J,\Sigma_J,\operatorname{Con}_J,\\
&\rho_J,O_J,K_J,B_J,\theta_J),
\qquad
\rho_J:X_J\to Y_J,\quad
O_J:X_J\to\mathcal O_J,\quad
\theta_J:X_J\to\{0,1\}.
\end{aligned}
\tag{#1}
\]
Its occurrence type is
\[
\begin{aligned}
\mathsf{Occ}_J:=\{&
(\mathfrak I_J;\Theta,x,\rho_J(x),\kappa,
\boldsymbol\varepsilon,\theta_J(x)):\\
&\Theta\in A_J,\ x\in X_J,\ \kappa\in K_J,\
\boldsymbol\varepsilon\in\RR_{\ge0}^{3}\}.
\end{aligned}
\tag{#2}
\]
An indexed occurrence $\mathsf{Occ}_{J,n}$ is an element of
$\mathsf{Occ}_J$.
}

\newcommand{\sharedmatching}{%
Let
\[
\mathcal A_8:=\{\infty\}\sqcup\ZZ_7,
\qquad
(\infty,0,1,2,3,4,5,6)
=(\mathsf{Gr},\mathsf{Li},\mathsf{Me},\mathsf{Ti},
\mathsf{Hz},\mathsf{Si},\mathsf{Sc},\mathsf{Mo}).
\]
For $r\in\ZZ_7$, set
\[
\mathsf M_r=\{\{\infty,r\}\}\cup
\{\{r+j,r-j\}:j=1,2,3\}.
\]
Then $E(K_8)=\bigsqcup_{r\in\ZZ_7}\mathsf M_r$, and the canonical matching is
\[
\mathsf M_0=\{\mathsf{Gravity\! -\! Light},\;
\mathsf{Memory\! -\! Motion},\;
\mathsf{Time\! -\! Scope},\;
\mathsf{Horizon\! -\! Silence}\}.
\]
}

\hypersetup{pdftitle={Gravity Light: Metric Invariance and Fixed Null Directions},pdfauthor={Anthony Vito Coppa},pdfsubject={Metric interval, null propagation, and local Lorentzian normal form}}
\begin{document}
\papertitle{Gravity Light}{Metric invariance and fixed null directions}

\begin{abstract}
Gravity measures interval. Light propagates along the null relation of that
interval. A recurrence matrix and an oriented area form generate a symmetric
split metric with two fixed null directions. Its metric flow is a Lorentzian
boost and its mirror reverses that boost. Every chosen Lorentzian two-plane of
spacetime has this root as an explicit normal form. Einstein curvature and
Maxwell curvature then describe the gravitational and electromagnetic fields
on the spacetime carrier. The resulting equations relate metric, null motion,
field strength, clock difference, and measured occurrence at one address.
\end{abstract}

\section{Metric field}

Physical interval is the Lorentzian quadratic form
\[
ds^2=g_{\mu\nu}\,dx^\mu dx^\nu.
\tag{GL.0}
\]
Gravity is the curvature relation of $g$. Light follows the null condition
\[
ds^2=0.
\tag{GL.0a}
\]
At each spacetime address $x\in M$, a Lorentzian two-plane
$P_x\subset T_xM$ carries one time direction, one transverse spatial
direction, and the two null directions between them. The following finite
root is an exact normal form of that local relation.

\section{The split metric}

\sharedroot

\begin{theorem}[Root identities]
The matrices above satisfy
\[
G=\begin{pmatrix}-2&1\\1&2\end{pmatrix},\qquad
G^2=5I,\qquad \det G=-5,\qquad
\operatorname{sig}(G)=(1,1).
\tag{GL.1}
\]
Moreover,
\[
S^{\mathsf T}G+GS=0,\qquad
H^{\mathsf T}GH=G,\qquad
HSH=-S.
\tag{GL.2}
\]
\end{theorem}

\begin{proof}
Direct multiplication gives $Q^2=Q+I$, $C^2=-I$, and
$G=Q\Omega-\Omega Q$. Put $A=2Q-I$. Then
\[
A^2=5I,\qquad G^2=5I,\qquad AG=5\Omega=-GA.
\]
Dividing by $\sqrt5$ gives the first two relations in GL.2. Since $H$ and
$S$ are symmetric involutions, the remaining identities follow directly.
\end{proof}

The metric flow is
\[
L_s=e^{sS}=\cosh(s)I+\sinh(s)S.
\tag{GL.3}
\]
Differentiating $L_s^{\mathsf T}GL_s$ and using GL.2 yields
\[
L_s^{\mathsf T}GL_s=G,\qquad
HL_sH=L_{-s}.
\tag{GL.4}
\]
Thus the mirror reverses the flow while preserving the metric relation.

\section{Null directions and fixed metric motion}

The null condition for $z=(x,y)^{\mathsf T}$ is
\[
z^{\mathsf T}Gz=-2x^2+2xy+2y^2=0.
\tag{GL.5}
\]
Its two projective solutions are
\[
\ell_\pm=\RR\left(1,\frac{-1\pm\sqrt5}{2}\right)^{\mathsf T}.
\tag{GL.6}
\]

\begin{proposition}[Fixed null directions]
Let
\[
A_G:=G\Omega=-\sqrt5S,\qquad
\mathsf B_t:=e^{tA_G}=L_{-\sqrt5t}.
\tag{GL.7}
\]
Then
\[
\mathsf B_t\ell_+=e^{-\sqrt5t}\ell_+,\qquad
\mathsf B_t\ell_-=e^{+\sqrt5t}\ell_-.
\tag{GL.8}
\]
\end{proposition}

\begin{proof}
The two lines are the eigendirections of $S$ with opposite eigenvalues. Apply
the exponential in GL.7.
\end{proof}

The two lines are the two null directions of the local split metric. The boost
changes scale along them in opposite directions while leaving each direction
fixed.

\section{Local Lorentzian geometry}

Define
\[
A_{\rm Lor}:=
\begin{pmatrix}
\sqrt{5/2}&0\\
1/\sqrt2&\sqrt2
\end{pmatrix},
\qquad
\eta_{1,1}:=\diag(-1,1).
\tag{GL.9}
\]
A direct multiplication gives
\[
A_{\rm Lor}^{\mathsf T}\eta_{1,1}A_{\rm Lor}=G.
\tag{GL.10}
\]
Consequently, with $\widetilde A=A_{\rm Lor}\oplus I_2$ and
$\eta_{3,1}=\diag(-1,1,1,1)$,
\[
\widetilde A^{\mathsf T}\eta_{3,1}\widetilde A=G\oplus I_2.
\tag{GL.11}
\]
This is an exact congruence of bilinear forms. On a curved Lorentzian
spacetime, choose at $x$ an isometric inclusion
\[
\iota_x:(\RR^2,G)\longrightarrow(P_x,g_x),
\qquad
g_x(\iota_xu,\iota_xv)=u^{\mathsf T}Gv.
\tag{GL.11a}
\]
The root null directions become physical local null directions:
\[
g_x(\iota_x\ell_\pm,\iota_x\ell_\pm)=0.
\tag{GL.11b}
\]
The root boost becomes the Lorentzian boost
\[
\mathsf B_t^{(x)}=\iota_x\mathsf B_t\iota_x^{-1}
\quad\text{on }P_x.
\tag{GL.11c}
\]

\section{Sevenfold transverse rotation}

Let $E=\RR^2$ carry the Euclidean metric $I_2$ and the oriented rotation
generator $C$. Define the four-dimensional carrier
\[
\widetilde V=\RR^2\oplus E,
\qquad
\widetilde G=G\oplus I_2,
\qquad
\widetilde C=C\oplus C.
\tag{GL.11d}
\]
With $D=\diag(1,-1)$, the mirror on the extended carrier is
\[
\widetilde H=H\oplus D.
\tag{GL.11e}
\]
The relation $DCD=-C$ combines with $HCH=-C$ to give
\[
\widetilde H^{\mathsf T}\widetilde G\widetilde H=\widetilde G,
\qquad
\widetilde H\widetilde C\widetilde H=-\widetilde C,
\qquad
\widetilde H^2=I_4.
\tag{GL.11f}
\]

For the cyclic generator $\tau$ of the seven-address transport, put
\[
R_r:=e^{(2\pi r/7)C},
\qquad
\varrho(\tau^r):=I_2\oplus R_r,
\qquad r\in\ZZ_7.
\tag{GL.11g}
\]
Then $\varrho$ is a faithful real representation of $\ZZ_7$ in the Lorentzian
isometry group:
\[
\varrho(\tau^r)\varrho(\tau^s)=\varrho(\tau^{r+s}),
\qquad
\varrho(\tau)^7=I_4,
\qquad
\varrho(\tau)^{\mathsf T}\widetilde G\varrho(\tau)=\widetilde G.
\tag{GL.11h}
\]
It commutes with phase and reverses under the mirror:
\[
\varrho(\tau)^{-1}\widetilde C\varrho(\tau)=\widetilde C,
\qquad
\widetilde H\varrho(\tau)\widetilde H=\varrho(\tau)^{-1}.
\tag{GL.11i}
\]
The embedded null directions are
\[
\widetilde\ell_\pm=(\ell_\pm,0)\in\widetilde V,
\qquad
\varrho(\tau^r)\widetilde\ell_\pm=\widetilde\ell_\pm.
\tag{GL.11j}
\]
The sevenfold transport rotates the transverse phase plane while leaving the
two local gravity--light null directions fixed. The mirror exchanges their
projective directions and reverses the cyclic rotation.

\section{Gravity and light}

The Einstein curvature tensor is
\[
\operatorname{Ein}_{\mu\nu}:=R_{\mu\nu}-\frac12R g_{\mu\nu}.
\tag{GL.12}
\]
The gravitational field equation is
\[
\operatorname{Ein}_{\mu\nu}+\Lambda g_{\mu\nu}
=\frac{8\pi G_N}{c^4}T_{\mu\nu}.
\tag{GL.12a}
\]
For an electromagnetic two-form $F$, Maxwell curvature obeys
\[
dF=0,\qquad d\star_gF=J.
\tag{GL.13}
\]
In vacuum $J=0$. In geometric optics, a wave covector $k$ is null when
\[
g^{-1}(k,k)=0.
\tag{GL.14}
\]

A plane-fronted gravity--light field has the form
\[
ds^2=2\,du\,dv+\mathcal H_{\rm pp}(u,x,y)\,du^2+dx^2+dy^2,
\qquad
R_{uu}=-\frac12\Delta_\perp\mathcal H_{\rm pp}.
\tag{GL.15}
\]
In vacuum, $R_{uu}=0$ gives $\Delta_\perp\mathcal H_{\rm pp}=0$. Thus curvature,
null propagation, and the local split metric meet in one spacetime field.

\section{Direct observation}

\sharedoccurrence{GL.16}{GL.17}

For $J=\mathsf{GL}$, take the state coordinates of $x_n$ to include the
local metric $g_n$, clock difference $\Delta\tau_n$, wave covector $k_n$, and
electromagnetic two-form $F_n$. With $O_n=O_{\mathsf{GL}}(x_n)$, the
gravity--light projection is
\[
\operatorname{pr}_{\rm GL}(\mathsf{Occ}_{\mathsf{GL},n})
=(\Theta_n,g_n,\Delta\tau_n,k_n,F_n,O_n,
\kappa_n,\boldsymbol\varepsilon_n).
\tag{GL.18}
\]
The physical relation is read by comparing these quantities at their common
address.

\end{document}
